The count is the dimension of $\mathfrak{so}(2\nu)$ on $2\nu$ Majorana generators (its elements are the bivectors, \claimref{std-algebra.bivector-lie-algebra}): the bivectors $\tfrac12[\gamma_a, \gamma_b]$ in the Majorana basis $\gamma_{2i-1} = v_i + v^*_i$, $\gamma_{2i} = J(v_i - v^*_i)$, with $J = e_2e_3$ the complex unit of the boundary modes (\claimref{synthesis.transverse-bivector-is-the-complex-unit-of-the-radial-atom}), are, up to scalars, the elements listed, the diagonal ones being $v^*_i v_i - v_i v^*_i = 2(v^*_i v_i - \tfrac12)$ by the anticommutation relations (\claimref{std-algebra.witt-basis-car}); without the $\tfrac12$ the bracket $[v^*_i v^*_j, v_j v_i] = v^*_i v_i + v^*_j v_j - 1$ leaves the span of the listed elements and lands in their span with the identity. On the vacuum $v_j|0\rangle = 0$ annihilates every off-diagonal hop and every annihilation, and $v^*_i v_i|0\rangle = 0$ gives the scalar $-\tfrac12$ for the Cartan elements; $v^*_i v^*_j|0\rangle = |ij\rangle$ are distinct basis states, hence independent and independent of $|0\rangle$, so the span of $X|0\rangle$ over $X \in \mathfrak{so}(2\nu)$ is that of $|0\rangle$ and the $|ij\rangle$, of dimension $1 + \nu(\nu - 1)/2$. The orbit of the ray of $|0\rangle$ under $\mathrm{Spin}(2\nu, \C)$ is the projective variety of pure spinors, of dimension $\nu(\nu - 1)/2$ (Cartan's theorem on pure spinors \cite{Chevalley1954}, \claimref{std-algebra.pure-spinors}), and the orbit of the vector is the cone over it, of dimension one more, so the creations span the whole tangent space of the ray orbit and, with the scalar direction, the whole tangent space of the vector orbit. The rotations are the exponentials of the corresponding generators, and the hop rotations fix $|0\rangle$ because their generators annihilate it. The script performs the count and the closure of the corrected generators under the bracket at $\nu = 2, 3, 4$, the annihilation and scalar statements, the tangent ranks $1 + \nu(\nu - 1)/2$ and $\nu(\nu - 1)/2$ by exact rank computations on the images of all $\nu(2\nu - 1)$ generators, the fixing and moving at $\nu = 4$ for every pair by the exponentials themselves, formed in closed form and verified orthogonal, the identification of the witness, and, as the converse the statement needs, that a hop rotation applied off the vacuum does entangle. A differential term of Dirac type, being first order in the modes, is a hop in this decomposition and lies in the stabiliser: of Remark \ref{rem:hopresidual}'s two residual types, only the pairing term is in the tangent.
